\section{Results and Discussion}
\subsection{Classification quality}

On the official AG News test split, the reported run evaluates 7,600
examples and obtains 93.7632\% accuracy, 93.7617\% macro F1, and a test
loss of 0.192954. The test set is balanced, with 1,900 examples in each
category. Sports has the highest class F1 (0.9851); Business has the
lowest (0.9047). The prediction file records 7,126 correct predictions
and 474 errors. The largest mutual confusion is between Business and
Sci/Tech: 139 Business examples are labeled Sci/Tech and 108 Sci/Tech
examples are labeled Business. These per-class results show why aggregate
accuracy alone does not describe all error patterns.

\begin{table}[htbp]
\centering
\small
\caption{Per-category classification results on the AG News test set.}
\label{tab:classification}
\begin{tabular}{lrrrr}
\toprule
Category & Precision & Recall & F1-score & Support \\
\midrule
World & 0.9552 & 0.9421 & 0.9486 & 1,900 \\
Sports & 0.9817 & 0.9884 & 0.9851 & 1,900 \\
Business & 0.9123 & 0.8974 & 0.9047 & 1,900 \\
Sci/Tech & 0.9017 & 0.9226 & 0.9121 & 1,900 \\
Macro average & 0.9377 & 0.9376 & 0.9376 & 7,600 \\
\bottomrule
\end{tabular}
\end{table}

The supplied result files do not identify the LoRA rank and do not
include a corresponding full fine-tuning result or results across ranks.
Accordingly, these metrics characterize the evaluated run only. They do
not show that LoRA matches or outperforms full fine-tuning, nor how rank
affects quality.

\subsection{Parameter and memory trade-off}

For one $1024\times1024$ matrix, full fine-tuning updates 1,048,576
parameters. LoRA trains $r(1024+1024)=2048r$ parameters: 8,192 at rank 4,
16,384 at rank 8, and 32,768 at rank 16. These correspond to reductions
of 99.21875\%, 98.4375\%, and 96.875\%, respectively, relative to the
single matrix. This illustrates a reduction in trainable parameters and
the associated gradients and optimizer states, but is not a count for
the full DistilRoBERTa model.

Trainable-parameter count is not total GPU memory use. The frozen base
weights remain in memory, and activations and temporary tensors also
contribute to runtime use. The supplied result files do not report peak
GPU memory, so they cannot establish an empirical memory saving.

\subsection{Training and inference efficiency}

The metrics report an evaluation runtime of 30.0006 seconds, or 253.328
test examples per second (15.833 evaluation steps per second). These are
evaluation measurements, not training-time measurements or a controlled
inference comparison. No training duration, peak memory, per-example
latency, or separate merged and unmerged adapter timings are reported.
Thus, efficiency benefits beyond the parameter-count calculation remain
method-level expectations rather than measured outcomes of this run.

LoRA's update can be merged into the base matrix after training, allowing
inference without a separate adapter branch. Keeping adapters separate
supports task switching, while merging makes the adapted weights easier
to use as a single model. O-LoRA can similarly accumulate task updates
for inference and retain individual factors for continued learning; that
choice entails growing adapter storage across tasks.

\subsection{Limitations}

The conclusions are limited by the single supplied evaluation, the
absence of rank-specific metadata and a full fine-tuning baseline, and
the lack of measured training time and GPU memory. No repeated runs or
variation estimates are available. The $1024\times1024$ example is an
illustration for one matrix, not a whole-model resource measurement.
Future controlled comparisons should record rank and target modules,
compare LoRA and full fine-tuning on the same split and schedule, and
measure training duration, peak GPU memory, and inference latency. O-LoRA
is discussed from its method description only and was not evaluated in
this experiment.

\subsection{When LoRA may not be the best choice}

LoRA is not automatically the right choice for every adaptation task. If
the task needs broad changes to model behavior, full fine-tuning may offer
more adaptation capacity. If the chosen rank cannot meet the quality
target, a higher rank or another method may be needed, with a corresponding
increase in trainable parameters. For small models with ample compute, the
efficiency gain may matter less than the simplicity of full fine-tuning.
Deployment requirements also matter: separate adapters ease task
switching, while merged weights simplify serving a selected adapted model.
The appropriate choice should be made against task-specific quality and
resource requirements.
